% ================= 2. PROJECT ORGANIZATION =================
\section{Project Organization}

\subsection{Team Organization}

The project is developed by a multidisciplinary team, with responsibilities
distributed across the main technical domains of the Cubli: mathematical
modeling, control, mechanical/CAD design, and system integration. Table
\ref{tab:team_roles} summarizes each member's primary role, and the detailed
responsibilities are outlined below.

\begin{table}[h]
\centering
\caption[Team members, contact details and primary roles]{Team members, contact details and primary roles.}
\label{tab:team_roles}
\small
\begin{tabular}{@{}p{3.6cm}p{4.6cm}p{6.0cm}@{}}
\toprule
\textbf{Member} & \textbf{Contact} & \textbf{Primary Role} \\
\midrule
Pablo Urioste Alarcón &
\href{mailto:pb.urioste4@gmail.com}{pb.urioste4@gmail.com}\newline +34 654 491 174 &
System integration and lead: digital actuator and sensor coding, embedded
software, and integration of the overall design. \\
\addlinespace
Deyan Nikolaev Vlaev &
\href{mailto:vlaev.dido@gmail.com}{vlaev.dido@gmail.com}\newline +359 877 082 499 &
Mathematical demonstration and derivation of the system dynamics. \\
\addlinespace
Niccolo &
--- &
Control system design and analysis. \\
\addlinespace
Andrea Liang &
\href{mailto:andrea.liang.2006@gmail.com}{andrea.liang.2006@gmail.com}\newline +31 06 5796 2976 &
Support on the control problem and structural/mechanical integration. \\
\addlinespace
Suvanna Viriyanti Wu &
\href{mailto:Viriyanti.wu@gmail.com}{Viriyanti.wu@gmail.com}\newline +39 329 074 7550 &
CAD modeling and mechanical design. \\
\addlinespace
Athanasia Nikolova &
--- &
CAD modeling, prototype building/assembly, and product presentation
(commercial pitch). \\
\bottomrule
\end{tabular}
\end{table}

\paragraph{Responsibilities.}
\begin{itemize}
    \item \textbf{Pablo Urioste Alarcón} --- System integration and overall
    coordination. Handles the coding of the digital actuators and sensors,
    the embedded software stack, and the integration of the mechanical,
    electronic and control subsystems into a working prototype.
    \item \textbf{Deyan Nikolaev Vlaev} --- Mathematical foundation. Responsible
    for the analytical demonstration and derivation of the equations of motion
    and the dynamic model of the Cubli.
    \item \textbf{Niccolo} --- Control. Leads the design, tuning and analysis of
    the control strategy used to balance and maneuver the Cubli.
    \item \textbf{Andrea Liang} --- Control and mechanics support. Assists on
    the control problem and contributes to the structural and mechanical
    integration of the system.
    \item \textbf{Suvanna Viriyanti Wu} --- CAD and mechanical design.
    Develops the 3D models and mechanical drawings of the structure and
    reaction-wheel assembly.
    \item \textbf{Athanasia Nikolova} --- CAD, prototype building and
    commercialization. Contributes to the 3D modeling, leads the physical
    building and assembly of the prototype, and is responsible for presenting
    and pitching the product.
\end{itemize}

\subsection{Work Breakdown Structure}
\label{sec:wbs}

This plan runs from project start to the final presentation on 20~August 2026.
The work is grouped into five streams --- investigation, requirements and
documentation; control and software; hardware design and build; three-dimensional
design; and electrical and electronics --- led by the owners assigned in
Table~\ref{tab:team_roles}. The investigation is kept to the essentials needed to
substantiate the design: a focused literature review, the dynamics derivation,
and the wheel sizing that follows from it. Table~\ref{tab:wbs_tasks} breaks the
streams into individual tasks, each with its scope, owner, dates and current
status.

The schedule was re-baselined on 31~July~2026, at the end of the second week, and
Table~\ref{tab:wbs_tasks} reflects that revision. The status column records the
position of each task at the re-baseline date: \emph{Complete} for work delivered
as planned, \emph{In progress} for work begun and continuing, \emph{Deferred} for
work not started and re-planned from the re-baseline date, \emph{New} for tasks
introduced by the revision, and \emph{Planned} for work whose start still lies
ahead. Tasks completed as planned retain their original dates, so the record of
what was undertaken in the first two weeks is preserved alongside the revised
plan rather than replaced by it.

% !TEX root = main.tex
% Cubli WBS - detailed task breakdown, 23 Jul - 20 Aug 2026.
%
% RE-BASELINED 31 July 2026 (project day 9).
% The original plan is preserved: tasks completed as planned retain their
% original dates and are marked "Complete". Tasks that were started but not
% finished are marked "In progress" and their dates are extended from the
% original start through to the re-baseline date and into the new plan. Tasks
% not begun are marked "Deferred" and re-planned from 31 July onward. Tasks
% introduced at the re-baseline are marked "New". No historical row has been
% deleted.
%
% Requires in preamble: longtable, booktabs, array, ragged2e + \newcolumntype{L}.

{\small
\setlength{\tabcolsep}{4pt}%
\begin{longtable}{@{}L{0.9cm} L{2.9cm} L{5.3cm} L{2.0cm} L{1.7cm} L{1.5cm}@{}}
\caption[Work breakdown structure: detailed task list, re-baselined 31 July 2026]{Work
breakdown structure, broken into individual tasks with scope, owner, dates and
status, grouped by workstream. Dates and status reflect the re-baseline of
31~July~2026; the original plan is retained for tasks completed as planned.}
\label{tab:wbs_tasks}\\
\toprule
\textbf{ID} & \textbf{Task} & \textbf{Scope} & \textbf{Owner} & \textbf{Dates} & \textbf{Status} \\
\midrule\endfirsthead
\toprule
\textbf{ID} & \textbf{Task} & \textbf{Scope} & \textbf{Owner} & \textbf{Dates} & \textbf{Status} \\
\midrule\endhead
\bottomrule\endfoot
\multicolumn{6}{@{}l}{\textbf{Investigation, Requirements \& Documentation}} \\
\midrule
A1 & Literature \& state of the art & Review the ETH Cubli, the one-wheel Cubli, CubeSat ADCS practice and the hopping-mission precedents & Suvanna \textit{(+Deyan)} & 23--25 Jul & Complete \\
A2 & Problem statement \& framing & Lock the ADCS application narrative that the rest of the document refers back to & Deyan & 23--24 Jul & Complete \\
A3 & Scope \& requirements & Define scope and build the requirements table (functional, performance, design constraint) with verification methods & Deyan \textit{(+Pablo)} & 24--25 Jul & Complete \\
A4 & Milestones \& roles & Prototype staging, binary gate criteria, and a one-page RACI of owners & Deyan & 23--25 Jul & Complete \\
A5 & Risk register & Trigger, impact, owner and mitigation for the leading risks & All \textit{(+Pablo)} & 28--29 Jul & Complete \\
A6 & 1-DoF equations of motion & Lagrangian derivation of the reaction-wheel pendulum with friction terms & Deyan \textit{(+Andrea)} & 24--25 Jul & Complete \\
A7 & 3D equations of motion & Coupled three-wheel dynamics for the edge and corner pivots & Deyan \textit{(+Andrea)} & 27--28 Jul & Complete \\
A8 & Linearisation \& controllability & Equilibria, linear models, controllability/observability checks & Deyan \textit{(+Andrea)} & 29 Jul & Complete \\
A9 & Jump-up sizing \& budget & Energy-method wheel sizing and the parametric mass/inertia/CoM budget & Deyan \textit{(+Andrea)} & 30 Jul--1 Aug & In progress \\
A10 & TDD draft (all sections) & Every member drafts their section; placeholder figures acceptable & All & 25--26 Jul & Complete \\
A11 & TDD deepen + rig evidence & Full modelling and control sections; insert the rig evidence and measured parameters available at the time of writing & All & 1--2 Aug & In progress \\
A12 & TDD v1 assemble \& submit & Consistency pass, figures, references, submission of the first assessed deliverable & All \textit{(Deyan integrates)} & 3 Aug & Planned \\
A13 & Assessor feedback triage & Receive the assessment of the v1 submission, triage the comments and assign each to an owner & Deyan \textit{(+All)} & 4--7 Aug & New \\
A14 & TDD v2 corrections \& second draft & Address the assessor comments and incorporate the build, bring-up and test results obtained since the v1 submission & All & 8--13 Aug & New \\
A15 & TDD v2 final submission & Final consistency pass and submission of the corrected document & All \textit{(Deyan integrates)} & 14 Aug & New \\
\midrule
\multicolumn{6}{@{}l}{\textbf{Control \& Software}} \\
\midrule
B1 & Sim environment + plant & Simulation setup and the 1-DoF plant, validated by energy conservation & Andrea & 1--3 Aug & Deferred \\
B2 & LQR design (1-DoF, sim) & Tune the regulator and sweep robustness to inertia error, delay and saturation & Andrea \textit{(+Deyan)} & 3--6 Aug & Deferred \\
B3 & State estimator & Complementary tilt filter with injected noise and bias in simulation & Nicc & 4--7 Aug & Deferred \\
B4 & Firmware \& comms & Firmware architecture, loop rates, and the CAN-FD command/reply layer & Nicc \textit{(+Pablo)} & 3--7 Aug & Deferred \\
B5 & Sensor drivers + cal & IMU and encoder drivers, gyro/accel calibration and the lever-arm correction & Nicc \textit{(+Andrea)} & 4--7 Aug & Deferred \\
B6 & Safety \& mode logic & Watchdog, tilt-limit disarm, wheel-speed cap and the mode state machine & Nicc \textit{(+Andrea)} & 5--7 Aug & Deferred \\
B7 & 1-DoF control design complete & Consolidated 1-DoF control stack --- plant, regulator, estimator and safety logic --- released as the baseline for the rig & Andrea, Nicc & 8--9 Aug & New \\
B8 & System identification & Measure $K_t$, wheel inertia, friction and loop latency on the rig & Andrea \textit{(+Nicc)} & 7--9 Aug & Planned \\
B9 & 3D model + LQR (sim) & 3D plant and the edge/corner regulators with the CAD inertia tensor & Andrea \textit{(+Deyan)} & 10--14 Aug & Planned \\
B10 & EKF / MEKF estimator & 3D estimator with gyro-bias states; characterise the yaw drift & Nicc \textit{(+Andrea)} & 12--15 Aug & Planned \\
B11 & Edge balance (S2a) & Retune the edge controller on the assembled cube; log disturbance recovery & Andrea \textit{(+Nicc)} & 15--17 Aug & Planned \\
B12 & Corner balance + slew (S2b/c) & Coupled three-axis corner balance and a commanded slew; report ADCS metrics & Andrea \textit{(+Nicc)} & 17--19 Aug & Planned \\
B13 & Jump-up (S3) & Spin-up profile and brake release for a face-to-edge jump. Sacrificial scope: yields to G6 and G7 if the schedule is contested & Nicc \textit{(+Andrea)} & 19--20 Aug & Planned \\
\midrule
\multicolumn{6}{@{}l}{\textbf{Hardware Design \& Build}} \\
\midrule
C1 & Envelope \& layout study & Packaging of wheels, drivers, battery and electronics with a central CoM & Neisa \textit{(+Suvanna)} & 23--24 Jul & Complete \\
C2 & Reaction-wheel design & Wheel to the inertia target: radius, rim, hub bolting to the motor bell & Suvanna \textit{(+Neisa)} & 24--26 Jul & Complete \\
C3 & Mounts \& brackets & Motor mount, ring-magnet carrier and adjustable encoder bracket & Neisa, Suvanna & 27--29 Jul & Complete \\
C4 & 1-DoF rig CAD + jigs & Single-axis rig and the balancing/assembly tooling & Suvanna \textit{(+Neisa)} & 25--28 Jul & Complete \\
C6 & Procurement order (round 1) & Order the long-lead items (rims, frame stock, connectors, spares) & Pablo \textit{(+Neisa)} & 23--24 Jul & Complete \\
C7 & Bench PSU \& bring-up & Current-limited supply, safety kit, and single-axis electrical bring-up & Pablo \textit{(+Nicc)} & 25--27 Jul & Complete \\
C8 & Fabricate wheel rings & Print the PET-CF rings and hubs. Ballasting is held at C8b pending the M6 fastener delivery & Neisa, Suvanna & 27--30 Jul & In progress \\
C8b & Wheel ballast \& completion & Install the M6 bolt/nut ballast stations to the sizing specification once the fasteners arrive, and complete each wheel & Suvanna \textit{(+Neisa)} & 4--6 Aug & New \\
C9 & Balance \& spin-test & Static-balance and containment spin-test every wheel, including a per-wheel ballast-count verification & Suvanna \textit{(+Neisa)} & 5--6 Aug & Deferred \\
C10 & Assemble 1-DoF rig & Motor, balanced wheel, encoder gap, IMU and hard stops. Includes the reinforced rig supports added at the re-baseline & Pablo \textit{(+Neisa)} & 1--2 Aug & In progress \\
\midrule
\multicolumn{6}{@{}l}{\textbf{3D Design}} \\
\midrule
D3a & Concept sizing \& fit study & Confirm that every component fits the \qty{150}{\milli\metre} envelope and lay out the internal structure & Neisa \textit{(+Suvanna)} & 2 Aug & New \\
D3b & Frame design freeze & Freeze the inertia-bearing frame geometry. Later features that do not affect the mass properties may still be added & Neisa, Suvanna & 2 Aug & New \\
D3c & Mass properties \& budget closure & Export the inertia tensor and CoM and close the mass budget against the sizing analysis & Neisa \textit{(+Andrea)} & 3--4 Aug & New \\
D3d & Detail CAD & Brackets, electronics bay, harness routing and the electrical keep-outs to the frozen board outline & Suvanna \textit{(+Neisa)} & 3--5 Aug & New \\
D3e & Procurement round 2 & Consolidate the bill of materials arising from the 3D design and place the second order & Pablo \textit{(+Neisa)} & 5 Aug & New \\
D3f & Tolerance \& clearance analysis & Fit and tolerance study, including the wheel sweep envelope and fastener access & Neisa \textit{(+Suvanna)} & 5--6 Aug & New \\
D3g & Manufacturing package release & Issue the drawings, print set and assembly instructions for fabrication & Neisa \textit{(+Suvanna)} & 7 Aug & New \\
C11 & Print structure + frame & Structural print set and the fabricated, orthogonal frame & Neisa, Suvanna & 8--12 Aug & Planned \\
C12 & Perfboard + harness & Power/signal board and the full power and CAN harness & Pablo & 8--12 Aug & Planned \\
C13 & Cube integration + power-on & Assemble the three axes, then bench and battery smoke test & Pablo \textit{(+Neisa)} & 13--15 Aug & Planned \\
C14 & Brake mechanism build & Fabricate and fit the servo brake; bench-test the deceleration & Suvanna \textit{(+Pablo)} & 15--18 Aug & Planned \\
\midrule
\multicolumn{6}{@{}l}{\textbf{Electrical \& Electronics}} \\
\midrule
E1 & Wiring \& placement design & Electrical architecture, wiring scheme and component placement for the single-axis rig, issued as the build reference & Pablo & 28--31 Jul & Complete \\
E2 & Rail, driver \& encoder bring-up & \qty{5}{\volt} rail soldered and trimmed to \qty{5.00}{\volt} under load; driver phase-wired and FOC-calibrated; MA600 encoder mounted at ${\le}\,\qty{0.5}{\milli\metre}$ air gap with verified read-back. Open-loop operation on the bench supply, no control loop & Pablo & 2 Aug & Planned \\
E2b & Bring-up recovery buffer & Reserved contingency to complete any part of E2 that does not close on the first attempt, protecting all downstream dates & Pablo & 3 Aug & New \\
E3 & Rig bus \& sensor integration & Two-node CAN-FD bus with split termination at both physical ends, Teensy and IMU integrated on the rig harness & Pablo \textit{(+Nicc)} & 4--5 Aug & Planned \\
E4 & Board outline \& floorplan & Driver placement from the encoder cable limit, CAN routing, and the perfboard floorplan issued as a frozen outline to CAD & Pablo \textit{(+Nicc)} & 4--5 Aug & Planned \\
E5 & 1-DoF closed-loop bring-up & Closed loop on the rig with the disarm path and watchdog verified before any balance attempt & Pablo \textit{(+Nicc)} & 6--7 Aug & Planned \\
E6 & Rig parameter measurement & Instrumented measurement of $K_t$, wheel inertia, friction and loop latency, delivered to the control workstream & Pablo \textit{(+Andrea)} & 7--9 Aug & Planned \\
E7 & Perfboard + battery power chain & Populate the distribution perfboard (two ground domains, one star point, split \qty{5}{\volt} rail); XT90 chain with anti-spark, fuse and ${\ge}\,\qty{50}{\volt}$ bulk capacitance; replicate the remaining actuation sets & Pablo \textit{(+Nicc)} & 8--11 Aug & Planned \\
E8 & 4-node bus + cube harness & Bus scaled to four nodes with stubs ${<}\,\qty{10}{\centi\metre}$ and an error soak at \qty{5}{\mega\bit\per\second}; full power and signal harness fabricated and routed; Wi-Fi telemetry for untethered operation & Pablo & 11--13 Aug & Planned \\
E9 & Cube power-on, \qty{48}{\ampere} + thermal & First battery power-on and the three-axis \qty{48}{\ampere} transient; driver thermals in the enclosed volume; all axes addressable & Pablo \textit{(+Nicc)} & 13--15 Aug & Planned \\
E10 & Brake servo rail + measurement & Servo on an isolated \qty{5}{\volt} rail with PWM drive and flyback; rail stability under stall, brake closure time and spin-down brake torque measured & Pablo \textit{(+Nicc)} & 16--19 Aug & Planned \\
\midrule
\multicolumn{6}{@{}l}{\textbf{Deliverables}} \\
\midrule
F1 & Results, demo \& final presentation & Slow-motion capture, results plots, deck, and the final presentation & All \textit{(+Suvanna edits)} & 18--20 Aug & Planned \\
\end{longtable}
}


\subsection{Schedule and Contingency Plan}
\label{sec:schedule}

The full-page schedule in Figure~\ref{fig:gantt} places every task of
Section~\ref{sec:wbs} on the 23~July--20~August 2026 calendar, and
Table~\ref{tab:gates} states the gate that closes each phase.

\paragraph{Position at the re-baseline.}
At the review of 31~July~2026 the investigation and documentation stream was
complete through the linearisation and controllability study, and the
single-axis mechanical and electrical designs were finished. Three areas were
behind the original plan. The control and software stream had not begun
implementation, the electrical work had produced the wiring and placement design
but no bench testing, and the three-dimensional design had not passed the
concept stage. The reaction-wheel rings were printed but could not be ballasted
or balanced, since the M6 fasteners that form the ballast stations had not yet
arrived. The revision below carries these items forward rather than compressing
them, and the completed work retains its original dates.

\paragraph{Structure of the revised plan.}
Sunday 2~August is the pivot of the schedule and carries two independent
objectives: the first powered bring-up of the single-axis rig, taking the
\qty{5}{\volt} rail, the driver and the encoder to verified open-loop operation
(M1); and the freeze of the three-dimensional frame geometry, which closes the
mass and inertia budget against the sizing analysis (M2). The following day is
held as a recovery buffer for the bring-up, so that a first-attempt failure on
the bench does not propagate into the dates that follow. The frame freeze
releases the second procurement round on 5~August (M3), which is the last point
at which parts can be ordered with useful delivery margin before assembly begins
on 13~August.

The documentation stream now carries two assessed submissions. The first is
delivered on 3~August (G4) with the evidence available at that date; the
assessor's comments are triaged during the following week and the corrected
second draft is submitted on 14~August (M8), by which point the rig results and
the integrated cube can both be reported.

\paragraph{Critical path and contingency.}
The critical path runs from the 2~August bring-up through the closed-loop rig
(M4), the measured plant parameters (M6) and the single-axis control design
(M7), into the three-dimensional controller and the balancing demonstrations.
The single-axis rig remains the protected item, since its balance result is the
experimental evidence in the document. The three-dimensional design proceeds in
parallel on the mechanical side, so that fabrication is not held by the control
work.

The plan retains the full scope through to the jump-up demonstration. The
schedule is nonetheless tightly constrained in the final week: edge balance
begins two days after cube assembly completes, and the jump-up gate falls on the
day of the final presentation. The jump-up (S3) is therefore designated
sacrificial scope by standing decision. If the balancing demonstrations require
the time, the jump-up yields first, and the demonstration presented on 20~August
is the corner balance and commanded slew. The brake mechanism build and its
associated electrical work are sequenced after the edge-balance gate for the
same reason.

\begin{table}[h]
\centering
\caption[Milestone gates, re-baselined 31 July 2026]{Milestone gates, from
kick-off to the final presentation, as revised at the re-baseline of
31~July~2026. The original date is retained for comparison; a dash indicates a
gate introduced by the revision.}
\label{tab:gates}
\small
\begin{tabular}{@{}l L{6.6cm} l l@{}}
\toprule
\textbf{Gate} & \textbf{Pass criterion} & \textbf{Original} & \textbf{Revised} \\
\midrule
G1 & Long-lead order placed; requirements and plan locked & 24 Jul & 24 Jul \\
G2 & Dynamics, linearisation and wheel sizing complete & 1 Aug & 31 Jul \\
M1 & 1-DoF rig bring-up: \qty{5}{\volt} rail trimmed under load, driver commutating, encoder verified at ${\le}\,\qty{0.5}{\milli\metre}$ gap & --- & 2 Aug \\
M2 & 3D frame geometry frozen at \qty{150}{\milli\metre}; mass, inertia and CoM budget closed against the sizing analysis & --- & 2 Aug \\
G4 & TDD v1 delivered for assessment & 3 Aug & 3 Aug \\
M3 & Procurement round 2 consolidated and ordered & --- & 5 Aug \\
M4 & 1-DoF closed loop live; disarm path and watchdog verified & --- & 7 Aug \\
M5 & Manufacturing package released for fabrication & --- & 7 Aug \\
M6 & Plant parameters measured on the rig and validated against simulation & --- & 9 Aug \\
M7 & 1-DoF control design complete: plant, regulator, estimator and safety logic & --- & 9 Aug \\
G3 & 1-DoF rig balances $\geq$60 s unassisted & 2 Aug & 12 Aug \\
M8 & TDD v2 corrected and finally submitted & --- & 14 Aug \\
G5 & Cube integrated and powered; all axes addressable & 9 Aug & 15 Aug \\
G6 & Edge balance (S2a) for $\geq$60 s with disturbance recovery & 13 Aug & 17 Aug \\
G7 & Corner balance and commanded slew (S2b/c) & 16 Aug & 19 Aug \\
G8 & Face-to-edge jump-up (S3) captured on video (sacrificial scope) & 18 Aug & 20 Aug \\
G9 & Final presentation and demonstration & 20 Aug & 20 Aug \\
\bottomrule
\end{tabular}
\end{table}

% !TEX root = ../main.tex
% ---------------------------------------------------------
% Cubli - consolidated project schedule, re-baselined 31 July 2026 (day 9).
%
% The chart merges the original plan with the re-baselined one on a single
% calendar. Tasks completed as planned keep their original bars in full
% saturation. Tasks that were started but not finished carry a pale bar over
% the original planned span, continued by a solid bar over the re-baselined
% span. Tasks not begun show the original span as a pale bar and the
% re-planned work as a solid bar. Nothing has been removed from the history.
%
% Day numbering is unchanged: day 1 = 23 Jul 2026, day 29 = 20 Aug 2026.
% The vertical rule at day 9 marks the re-baseline date.
% ---------------------------------------------------------

\ganttset{
  hgrid, vgrid={*6{draw=none}, *1{draw=black!15}},
  x unit=0.4cm, y unit title=0.5cm, y unit chart=0.315cm,
  title height=1, bar height=0.6, bar top shift=0.2,
  group top shift=0.1, group height=0.5, milestone top shift=0.15,
  bar label font=\small, group label font=\bfseries\small,
  milestone label font=\small\itshape, title label font=\small,
  bar/.append style={draw=none},
}

\begin{figure}[p]
\centering
\resizebox{\textwidth}{!}{%
\begin{ganttchart}[
    today=9, today rule/.style={draw=black!70, dashed, line width=0.7pt},
    today label={}
  ]{1}{29}
  \gantttitle{Week 1 (23--29 Jul)}{7}
  \gantttitle{Week 2 (30 Jul--5 Aug)}{7}
  \gantttitle{Week 3 (6--12 Aug)}{7}
  \gantttitle{Week 4 (13--19 Aug)}{7}
  \gantttitle{20}{1} \\
  \gantttitle{23}{1} \gantttitle{24}{1} \gantttitle{25}{1} \gantttitle{26}{1} \gantttitle{27}{1} \gantttitle{28}{1} \gantttitle{29}{1} \gantttitle{30}{1} \gantttitle{31}{1} \gantttitle{1}{1} \gantttitle{2}{1} \gantttitle{3}{1} \gantttitle{4}{1} \gantttitle{5}{1} \gantttitle{6}{1} \gantttitle{7}{1} \gantttitle{8}{1} \gantttitle{9}{1} \gantttitle{10}{1} \gantttitle{11}{1} \gantttitle{12}{1} \gantttitle{13}{1} \gantttitle{14}{1} \gantttitle{15}{1} \gantttitle{16}{1} \gantttitle{17}{1} \gantttitle{18}{1} \gantttitle{19}{1} \gantttitle{20}{1} \\
  % ===== Investigation, Requirements & Documentation =====
  \ganttgroup[group/.append style={fill=pfcol!55}]{Documentation}{1}{23} \ganttnewline
  \ganttbar[bar/.append style={fill=pfcol}]{A1 Literature review}{1}{3} \ganttnewline
  \ganttbar[bar/.append style={fill=pfcol}]{A2 Problem statement}{1}{2} \ganttnewline
  \ganttbar[bar/.append style={fill=pfcol}]{A3 Scope \& requirements}{2}{3} \ganttnewline
  \ganttbar[bar/.append style={fill=pfcol}]{A4 Milestones \& roles}{1}{3} \ganttnewline
  \ganttbar[bar/.append style={fill=pfcol}]{A5 Risk register}{6}{7} \ganttnewline
  \ganttbar[bar/.append style={fill=pfcol}]{A6 1-DoF EoM}{2}{3} \ganttnewline
  \ganttbar[bar/.append style={fill=pfcol}]{A7 3D EoM}{5}{6} \ganttnewline
  \ganttbar[bar/.append style={fill=pfcol}]{A8 Linearisation}{7}{7} \ganttnewline
  \ganttbar[bar/.append style={fill=pfcol}]{A9 Jump-up sizing \& budget}{8}{10} \ganttnewline
  \ganttbar[bar/.append style={fill=pfcol}]{A10 TDD draft}{3}{4} \ganttnewline
  \ganttbar[bar/.append style={fill=pfcol}]{A11 TDD deepen + rig evidence}{10}{11} \ganttnewline
  \ganttbar[bar/.append style={fill=pfcol}]{A12 TDD v1 submit}{12}{12} \ganttnewline
  \ganttbar[bar/.append style={fill=pfcol}]{A13 Assessor feedback triage}{13}{16} \ganttnewline
  \ganttbar[bar/.append style={fill=pfcol}]{A14 TDD v2 corrections}{17}{22} \ganttnewline
  \ganttbar[bar/.append style={fill=pfcol}]{A15 TDD v2 final submission}{23}{23} \ganttnewline
  % ===== Control & Software =====
  \ganttgroup[group/.append style={fill=ctcol!55}]{Control \& Software}{10}{29} \ganttnewline
  \ganttbar[bar/.append style={fill=ctcol!25}]{B1 Sim env + plant (as planned)}{1}{2} \ganttnewline
  \ganttbar[bar/.append style={fill=ctcol}]{B1 Sim env + plant}{10}{12} \ganttnewline
  \ganttbar[bar/.append style={fill=ctcol!25}]{B2 LQR 1-DoF (as planned)}{4}{5} \ganttnewline
  \ganttbar[bar/.append style={fill=ctcol}]{B2 LQR 1-DoF (sim)}{12}{15} \ganttnewline
  \ganttbar[bar/.append style={fill=ctcol!25}]{B3 State estimator (as planned)}{5}{6} \ganttnewline
  \ganttbar[bar/.append style={fill=ctcol}]{B3 State estimator}{13}{16} \ganttnewline
  \ganttbar[bar/.append style={fill=ctcol!25}]{B4 Firmware \& comms (as planned)}{2}{6} \ganttnewline
  \ganttbar[bar/.append style={fill=ctcol}]{B4 Firmware \& comms}{12}{16} \ganttnewline
  \ganttbar[bar/.append style={fill=ctcol}]{B5 Sensor drivers + cal}{13}{16} \ganttnewline
  \ganttbar[bar/.append style={fill=ctcol}]{B6 Safety \& mode logic}{14}{16} \ganttnewline
  \ganttbar[bar/.append style={fill=ctcol}]{B7 1-DoF control design complete}{17}{18} \ganttnewline
  \ganttbar[bar/.append style={fill=ctcol}]{B8 System identification}{16}{18} \ganttnewline
  \ganttbar[bar/.append style={fill=ctcol}]{B9 3D model + LQR (sim)}{19}{23} \ganttnewline
  \ganttbar[bar/.append style={fill=ctcol}]{B10 EKF / MEKF estimator}{21}{24} \ganttnewline
  \ganttbar[bar/.append style={fill=ctcol}]{B11 Edge balance (S2a)}{24}{26} \ganttnewline
  \ganttbar[bar/.append style={fill=ctcol}]{B12 Corner balance + slew}{26}{28} \ganttnewline
  \ganttbar[bar/.append style={fill=ctcol}]{B13 Jump-up (S3)}{28}{29} \ganttnewline
  % ===== Hardware Design & Build =====
  \ganttgroup[group/.append style={fill=cdcol!55}]{Hardware Design \& Build}{1}{15} \ganttnewline
  \ganttbar[bar/.append style={fill=cdcol}]{C1 Envelope \& layout}{1}{2} \ganttnewline
  \ganttbar[bar/.append style={fill=cdcol}]{C2 Reaction-wheel design}{2}{4} \ganttnewline
  \ganttbar[bar/.append style={fill=cdcol}]{C3 Mounts \& brackets}{5}{7} \ganttnewline
  \ganttbar[bar/.append style={fill=cdcol}]{C4 1-DoF rig CAD + jigs}{3}{6} \ganttnewline
  \ganttbar[bar/.append style={fill=cdcol}]{C6 Procurement round 1}{1}{2} \ganttnewline
  \ganttbar[bar/.append style={fill=cdcol}]{C7 Bench PSU \& bring-up}{3}{5} \ganttnewline
  \ganttbar[bar/.append style={fill=cdcol}]{C8 Fabricate wheel rings}{5}{9} \ganttnewline
  \ganttbar[bar/.append style={fill=cdcol}]{C8b Wheel ballast \& completion}{13}{15} \ganttnewline
  \ganttbar[bar/.append style={fill=cdcol!25}]{C9 Balance \& spin-test (as planned)}{8}{9} \ganttnewline
  \ganttbar[bar/.append style={fill=cdcol}]{C9 Balance \& spin-test}{14}{15} \ganttnewline
  \ganttbar[bar/.append style={fill=cdcol}]{C10 Assemble 1-DoF rig}{10}{11} \ganttnewline
  % ===== 3D Design =====
  \ganttgroup[group/.append style={fill=cdcol!55}]{3D Design}{11}{27} \ganttnewline
  \ganttbar[bar/.append style={fill=cdcol}]{D3a Concept sizing \& fit study}{11}{11} \ganttnewline
  \ganttbar[bar/.append style={fill=cdcol}]{D3b Frame design freeze}{11}{11} \ganttnewline
  \ganttbar[bar/.append style={fill=cdcol}]{D3c Mass properties \& budget}{12}{13} \ganttnewline
  \ganttbar[bar/.append style={fill=cdcol}]{D3d Detail CAD}{12}{14} \ganttnewline
  \ganttbar[bar/.append style={fill=cdcol}]{D3e Procurement round 2}{14}{14} \ganttnewline
  \ganttbar[bar/.append style={fill=cdcol}]{D3f Tolerance \& clearance}{14}{15} \ganttnewline
  \ganttbar[bar/.append style={fill=cdcol}]{D3g Manufacturing package}{16}{16} \ganttnewline
  \ganttbar[bar/.append style={fill=cdcol}]{C11 Print structure + frame}{17}{21} \ganttnewline
  \ganttbar[bar/.append style={fill=cdcol}]{C12 Perfboard + harness}{17}{21} \ganttnewline
  \ganttbar[bar/.append style={fill=cdcol}]{C13 Cube integration + power-on}{22}{24} \ganttnewline
  \ganttbar[bar/.append style={fill=cdcol}]{C14 Brake mechanism build}{24}{27} \ganttnewline
  % ===== Electrical & Electronics =====
  \ganttgroup[group/.append style={fill=mfcol!55}]{Electrical \& Electronics}{6}{28} \ganttnewline
  \ganttbar[bar/.append style={fill=mfcol}]{E1 Wiring \& placement design}{6}{9} \ganttnewline
  \ganttbar[bar/.append style={fill=mfcol}]{E2 Rail, driver \& encoder bring-up}{11}{11} \ganttnewline
  \ganttbar[bar/.append style={fill=mfcol!45}]{E2b Bring-up recovery buffer}{12}{12} \ganttnewline
  \ganttbar[bar/.append style={fill=mfcol}]{E3 Rig bus \& sensor integration}{13}{14} \ganttnewline
  \ganttbar[bar/.append style={fill=mfcol}]{E4 Board outline \& floorplan}{13}{14} \ganttnewline
  \ganttbar[bar/.append style={fill=mfcol}]{E5 1-DoF closed-loop bring-up}{15}{16} \ganttnewline
  \ganttbar[bar/.append style={fill=mfcol}]{E6 Rig parameter measurement}{16}{18} \ganttnewline
  \ganttbar[bar/.append style={fill=mfcol}]{E7 Perfboard + battery chain}{17}{20} \ganttnewline
  \ganttbar[bar/.append style={fill=mfcol}]{E8 4-node bus + cube harness}{20}{22} \ganttnewline
  \ganttbar[bar/.append style={fill=mfcol}]{E9 Cube power-on, 48\,A + thermal}{22}{24} \ganttnewline
  \ganttbar[bar/.append style={fill=mfcol}]{E10 Brake servo rail + measurement}{25}{28} \ganttnewline
  % ===== Deliverables =====
  \ganttgroup[group/.append style={fill=mscol!45}]{Deliverables}{27}{29} \ganttnewline
  \ganttbar[bar/.append style={fill=mscol!70}]{F1 Results, demo \& presentation}{27}{29} \ganttnewline
  \ganttnewline
  % ===== Milestone gates =====
  \ganttmilestone[milestone/.append style={fill=mscol}]{G1 Kick-off}{2} \ganttnewline
  \ganttmilestone[milestone/.append style={fill=mscol}]{G2 Modelling complete}{9} \ganttnewline
  \ganttmilestone[milestone/.append style={fill=mscol}]{M1 1-DoF electrical bring-up}{11} \ganttnewline
  \ganttmilestone[milestone/.append style={fill=mscol}]{M2 3D frame design frozen}{11} \ganttnewline
  \ganttmilestone[milestone/.append style={fill=mscol}]{G4 TDD v1 delivered}{12} \ganttnewline
  \ganttmilestone[milestone/.append style={fill=mscol}]{M3 Procurement round 2 placed}{14} \ganttnewline
  \ganttmilestone[milestone/.append style={fill=mscol}]{M4 1-DoF closed loop live}{16} \ganttnewline
  \ganttmilestone[milestone/.append style={fill=mscol}]{M5 Manufacturing package released}{16} \ganttnewline
  \ganttmilestone[milestone/.append style={fill=mscol}]{M6 Plant parameters validated}{18} \ganttnewline
  \ganttmilestone[milestone/.append style={fill=mscol}]{M7 1-DoF control design complete}{18} \ganttnewline
  \ganttmilestone[milestone/.append style={fill=mscol}]{G3 1-DoF rig balances}{21} \ganttnewline
  \ganttmilestone[milestone/.append style={fill=mscol}]{M8 TDD v2 final submitted}{23} \ganttnewline
  \ganttmilestone[milestone/.append style={fill=mscol}]{G5 Cube integrated}{24} \ganttnewline
  \ganttmilestone[milestone/.append style={fill=mscol}]{G6 Edge balance}{26} \ganttnewline
  \ganttmilestone[milestone/.append style={fill=mscol}]{G7 Corner balance + slew}{28} \ganttnewline
  \ganttmilestone[milestone/.append style={fill=mscol}]{G8 Jump-up}{29} \ganttnewline
  \ganttmilestone[milestone/.append style={fill=mscol}]{G9 Final presentation}{29}
\end{ganttchart}%
}
\caption[Consolidated project schedule, re-baselined 31 July 2026]{Consolidated
project schedule, 23~July--20~August 2026, re-baselined on 31~July~2026 (dashed
vertical rule). Work completed as planned is shown in full saturation on its
original dates. Where a task was planned but not completed, the original span is
retained as a pale bar and the re-baselined work follows as a solid bar of the
same colour, so that the plan of record and the revised plan are both legible on
one calendar. The documentation stream (green) now carries two assessed
submissions, on 3~and 14~August. Control and software (blue) is re-planned from
the re-baseline date, with the three-dimensional extension concentrated in the
final two weeks. Hardware (gold) separates the single-axis build from the
three-dimensional design, and the electrical stream (orange) proceeds from the
single-axis bring-up on 2~August to the integrated cube. Grey diamonds mark the
milestone gates listed in Table~\ref{tab:gates}.}
\label{fig:gantt}
\end{figure}

